\documentclass[a4paper,12pt]{letter}
\usepackage[utf8]{inputenc}
\usepackage[english,slovene]{babel}
\usepackage[tikz]{bclogo}
\usepackage{qrcode,marvosym}
\usepackage{pgf,tikz,pgfplots}
\pgfplotsset{compat=1.14}
\usepackage{mathrsfs}
\usetikzlibrary{arrows}
\usepackage{graphicx}
\newcommand{\ds}{\displaystyle}
\usepackage{fancybox}
\usepackage{amsfonts}
\usepackage{amsmath}
\usepackage{wallpaper}
\usepackage{hyperref}
\usepackage[cp1250]{inputenc}
\usepackage{array}
\newcolumntype{M}[1]{>{\centering\arraybackslash}m{#1}}
\newcolumntype{N}{@{}m{0pt}@{}}
\usepackage{fancyhdr,enumerate}

\newcommand{\ora}{\overrightarrow}
\renewcommand{\headrulewidth}{3pt}
\renewcommand{\headrule}{\hbox to\headwidth{%
  \color{gray}\leaders\hrule height \headrulewidth\hfill}}
\renewcommand{\footrulewidth}{2pt}
\renewcommand{\footrule}{\hbox to\headwidth{%
  \color{brown}\leaders\hrule height \footrulewidth\hfill}}
\usepackage{gensymb}
\usepackage{amssymb}

\usepackage{fancyhdr}
\textwidth 20 true cm
\oddsidemargin -2 true cm
\evensidemargin -2 true cm
\topmargin -3.5 true cm
\textheight 27.5 true cm
\linespread{1.2}
\renewcommand{\headrulewidth}{0.3pt}
\renewcommand{\footrulewidth}{1.9pt}
\definecolor{qqwuqq}{rgb}{1,0,0}
\definecolor{qqqqff}{rgb}{0,0,1}
\definecolor{qqffqq}{rgb}{0,1,0}
\definecolor{ffqqqq}{rgb}{1,0,0}
\definecolor{zzuxux}{rgb}{0,0,0}
\usepackage{mdframed}
\usepackage{multicol}
\definecolor{bgblue}{RGB}{0,0,253}
\usepackage{xcolor}
\usepackage{eurosym}

\newcounter{tocke}
\setcounter{tocke}{0}
\usepackage{tikz-3dplot}

\mdfdefinestyle{theoremstyle}{%
linecolor=brown!40,linewidth=1pt,%
frametitlerule=true,%
frametitlebackgroundcolor=brown!50,
innertopmargin=\topskip,
}

\mdfdefinestyle{theoremstyle1}{%
linecolor=brown,middlelinewidth=1pt,%
frametitlerule=true,%
apptotikzsetting={\tikzset{mdfframetitlebackground/.append style={%
shade,left color=black!40, right color=gray!10},
mdfbackground/.append style={
top color=white!100!white,
bottom color=black!5
},
}},
frametitlerulecolor=gray,
frametitlerulewidth=2pt,
innertopmargin=\topskip,
innerbottommargin=\topskip,
}
\mdtheorem[style=theoremstyle1]{naloga}{Naloga}
\mdtheorem[style=theoremstyle]{resitev}{Analiza Naloge}
\mdtheorem[style=theoremstyle]{kriterij}{\v{S}tevilo dose\v{z}enih to\v{c}k na testu}

\usepackage[table]{xcolor}
\usepackage{titlesec}
\usepackage{venndiagram}
\usepackage[printwatermark]{xwatermark}
\usepackage{xcolor}
\newwatermark[allpages,color=brown!2,angle=45,scale=5,
xpos=-5,ypos=0]{matej.info}

\titleformat{\subsection}
  {\normalfont\Large\bfseries\color{brown}}
  {\thesubsection}{1em}{}

\titleformat{\subsubsection}
  {\normalfont\large\bfseries\color{brown}}
  {\thesubsubsection}{1em}{}

\begin{document}
\pagestyle{fancy}
\renewcommand\headrulewidth{0pt}
\lhead{}\chead{}\rhead{}
\rfoot{\vspace*{-2\baselineskip}}

\renewcommand{\vec}{\overrightarrow}
\everymath{\displaystyle}

\begin{minipage}{20cm}
\begin{bclogo}[couleur=brown!40,
  arrondi=0,ombre=true, couleurOmbre=gray!50,
  logo=\bcplume, noborder=true]
{\textrm{{\Large{Popravni izpit 2PTI} \Huge{}:
\large{2. letnik} \hfill $PTI-2$}}}
\end{bclogo}
\end{minipage}

\begin{minipage}{20cm}
\begin{bclogo}[couleur=brown!40,ombre=true, couleurOmbre=gray!50,
  arrondi=0, logo=\bcplume, noborder=true]
{\textsc{{\large{Ime in priimek:}
\vspace{0.2cm} \rule{15cm}{0.2mm}\hfill }}}
\end{bclogo}
\end{minipage}

%==============================================
% NALOGE
%==============================================

\def\tockeA{3+2+3+2}
\begin{naloga}[\hfill $\tockeA$] \addtocounter{tocke}{\tockeA}
Podan je graf polinoma tretje stopnje.

a) Določi ničle in njihovo večkratnost.

b) Določi $p(0)$ in $p(-1)$.

c) Zapiši predpis polinoma.

d) Kje velja $p(x)>0$?

\end{naloga}

\vspace{1cm}


\begin{tikzpicture}[line cap=round,line join=round,>=triangle 45,x=0.6cm,y=.6cm]
\begin{axis}[
x=0.6cm,y=0.6cm,
axis lines=middle,
ymajorgrids=true,
xmajorgrids=true,
xmin=-4.5,xmax=3.5,
ymin=-7,ymax=7,
xtick={-4,-3,...,3},
ytick={-6,-5,...,6},
width=10cm,height=8cm,
]
\addplot[line width=1.pt,color=qqwuqq,smooth,samples=100,domain=-3.6:3.0] {-(x-1)*(x+2)^2};
\node[color=qqwuqq] at (axis cs:-2.6,5) {$f$};
\end{axis}
\end{tikzpicture}
\newpage

\def\tockeB{7+3}
\begin{naloga}[\hfill $\tockeB$] \addtocounter{tocke}{\tockeB}
a) Izračunaj količnik in ostanek pri deljenju:
\[
p(x)=3x^3-2x^2+5x-1, \quad q(x)=x^2+x-1.
\]

b) Polinom $q(x)=(a+1)x^3+2bx+cx-1$ je enak polinomu $p(x)$. Določi $a,b,c.$
\end{naloga}

\newpage

\def\tockeC{3+3+4}
\begin{naloga}[\hfill $\tockeC$
\qquad \rightsquigarrow \vert a.\qquad|b.\qquad|c.\qquad|] \addtocounter{tocke}{\tockeC}
Reši enačbo:
\vspace{-0.5cm}
\begin{multicols}{3}
a) $\log_3(2x-1)=3$

b) $5^x=8^{-x}$

c) $3^x=20$
\end{multicols}
\end{naloga}

\newpage


\def\tockeD{5+5}
\begin{naloga}[\hfill $\tockeD$
\qquad \rightsquigarrow \vert a.\qquad|b.\qquad|] \addtocounter{tocke}{\tockeD}
V aritmetičnem zaporedju je prvi člen enak $12$, peti člen pa $-8$.

a) Izračunaj diferenco zaporedja in tretji člen.

b) Izračunaj vsoto prvih 12 členov tega zaporedja.
\end{naloga}

\newpage

\def\tockeE{4+3+3}
\begin{naloga}[\hfill $\tockeE$
\qquad \rightsquigarrow \vert a.\qquad|b.\qquad|c.\qquad|] \addtocounter{tocke}{\tockeE}
Pri metih kocke so bili doseženi naslednji rezultati:
\[
2,5,3,6,1,4,2,3,5,6,2,1,4,3,2,5,6,3,4,1,2,3,5,4.
\]

a) Izračunaj relativne frekvence za posamično vrednost.

b) Prikaži rezultate s histogramom.

c) Izračunaj povprečno vrednost, modus in mediano.
\end{naloga}

\newpage

\def\tockeF{5+5}
\begin{naloga}[\hfill $\tockeF$
\qquad \rightsquigarrow \vert a.\qquad|b.\qquad|] \addtocounter{tocke}{\tockeF}
V skupini je 6 fantov in 4 dekleta.

a) Na koliko načinov lahko izberemo odbor treh članov, če mora biti v odboru vsaj en fant?

b) Kolikšna je verjetnost, da bo v naključno izbranem odboru treh članov same dekleta?
\end{naloga}

\newpage

\def\tockeG{5+5}
\begin{naloga}[\hfill $\tockeG$
\qquad \rightsquigarrow \vert a.\qquad|b.\qquad|] \addtocounter{tocke}{\tockeG}
a) Izračunaj odvod funkcije $f(x)=4x^3-3x^2+2x$ v točki $x_0=1$

b) Kakšen je kot med funkcijo$g(x)=\dfrac{4}{\sqrt[3]{x^2}}-3$  in abscisno osjo  v $x_0=-1$?
\end{naloga}

\newpage

\def\tockeH{6+4}
\begin{naloga}[\hfill $\tockeH$
\qquad \rightsquigarrow \vert a.\qquad|b.\qquad|] \addtocounter{tocke}{\tockeH}
Določi enačbo tangente na graf funkcije $f(x)=x^2-1$ v točki $x_0=3$. Nariši funkcijo in tangento v koordinatni sistem.
\end{naloga}

\vfill
\mdtheorem[style=theoremstyle1]{kriterij}{Kriterij ocenjevanja}
\begin{kriterij*}[\hfill \v{s}tevilo vseh to\v{c}k na testu: $\arabic{tocke}$]
$$
\begin{array}{||c|c|c|c|c|c||>{}c|c|}
\hline\hline
\textup{ocena}& 1&2&3&4&5&\textrm{uspe\v{s}nost v } \%& \textrm{\bf OCENA}\\
\hline\hline
\% & [0,45) &[45,60)&[60,75) &[75,90) & [90,100]&
\mbox{\phantom{\huge{$\frac{5}{445}$15}}} & \\
\hline
\end{array}
\hspace*{0.5cm}\vspace*{-0.3cm}
\qrcode[hyperlink,height=0.8in]{http://www.matej.info}
$$
\vspace*{0.1cm}
\end{kriterij*}
\newpage


%==============================================
% REŠITVE
%==============================================

\begin{minipage}{20cm}
\begin{bclogo}[couleur=brown!40,
  arrondi=0,ombre=true, couleurOmbre=gray!50,
  logo=\bcplume, noborder=true]
{\textrm{{\Large{Popravni izpit 2PTI} \Huge{}:
\large{Re\v{s}itve} \hfill $PTI-2$}}}
\end{bclogo}
\end{minipage}

\vspace{0.5cm}

\begin{resitev}[Naloga 1 --- Graf polinoma]

Podan je graf:
\[
f(x)=-(x-1)(x+2)^2
\]

a) Ničle:
\[
x=1 \ (\text{enostavna}), \quad x=-2 \ (\text{dvojna})
\]

b)
\[
p(0)=-(0-1)(0+2)^2=-(-1)(4)=4
\]
\[
p(-1)=-(-1-1)(-1+2)^2=-(-2)(1)=2
\]

c)
\[
p(x)=-(x-1)(x+2)^2
\]

d) Predznak funkcije: $(x+2)^2\ge 0$ za vse $x$, zato je predznak $f(x)$ odvisen od $-(x-1)$.
\[
-(x-1)>0 \;\Leftrightarrow\; x<1
\]
V točki $x=-2$ je $f(x)=0$ (dvojna ničla), zato to točko izvzamemo. Torej:
\[
p(x)>0 \;\Leftrightarrow\; x\in(-\infty,-2)\cup(-2,1)
\]
\end{resitev}

\vspace{0.5cm}

\begin{resitev}[Naloga 2 --- Deljenje in primerjava polinomov]

a)
\[
\frac{3x^3-2x^2+5x-1}{x^2+x-1}
\]
S postopnim deljenjem dobimo:
\[
k(x)=3x-5,\qquad r(x)=13x-6
\]

b) Polinom $q(x)=(a+1)x^3+bx^2+cx-1$ mora biti enak polinomu
\[
p(x)=3x^3-2x^2+5x-1.
\]
Primerjamo ustrezne koeficiente:
\[
a+1=3 \;\Rightarrow\; a=2
\]
\[
b=-2
\]
\[
c=5
\]
\textit{(Opomba: pri primerjavi koeficientov polinomov \v{c}len $x^2$ ne sme manjkati, sicer enakost ni mo\v{z}na, saj ima $p(x)$ koeficient $-2$ pri $x^2$.)}
\end{resitev}

\newpage

\begin{resitev}[Naloga 3 --- Eksponentne in logaritemske enačbe]

a) $\log_3(2x-1)=3$
\[
2x-1=3^3=27 \;\Rightarrow\; 2x=28 \;\Rightarrow\; x=14
\]

b) $5^x=8^{-x}$
\[
5^x\cdot 8^x=1 \;\Rightarrow\; (5\cdot8)^x=1 \;\Rightarrow\; 40^x=1 \;\Rightarrow\; x=0
\]

c) $3^x=20$
\[
x=\log_3 20=\frac{\log 20}{\log 3}\approx 2{,}727
\]
\end{resitev}

\vspace{0.5cm}

\begin{resitev}[Naloga 4 --- Aritmetično zaporedje]

Podatki: $a_1=12,\; a_5=-8$.

a)
\[
a_5=a_1+4d \;\Rightarrow\; -8=12+4d \;\Rightarrow\; d=-5
\]
\[
a_3=a_1+2d=12-10=2
\]

b)
\[
S_{12}=\frac{12}{2}\big(2a_1+11d\big)=6\big(24-55\big)=6\cdot(-31)=-186
\]
\end{resitev}

\newpage

\begin{resitev}[Naloga 5 --- Statistika]

Podatki (24 vrednosti): frekvence po skupinah so:
\[
1:3,\quad 2:5,\quad 3:5,\quad 4:4,\quad 5:4,\quad 6:3
\]

a) Relativne frekvence:
\[
f_1=\tfrac{3}{24}\approx0{,}125,\; f_2=\tfrac{5}{24}\approx0{,}208,\; f_3=\tfrac{5}{24}\approx0{,}208,
\]
\[
f_4=\tfrac{4}{24}\approx0{,}167,\; f_5=\tfrac{4}{24}\approx0{,}167,\; f_6=\tfrac{3}{24}\approx0{,}125
\]

b) Histogram:
\begin{center}
\begin{tikzpicture}
\begin{axis}[
ybar, bar width=0.8cm,
xlabel={Vrednost},
ylabel={Frekvenca},
symbolic x coords={1,2,3,4,5,6},
xtick=data,
ymin=0, ymax=6,
width=10cm, height=6cm,
nodes near coords,
]
\addplot coordinates {(1,3) (2,5) (3,5) (4,4) (5,4) (6,3)};
\end{axis}
\end{tikzpicture}
\end{center}

c)
\[
\bar{x}=\frac{1\cdot3+2\cdot5+3\cdot5+4\cdot4+5\cdot4+6\cdot3}{24}=\frac{82}{24}\approx3{,}42
\]
\[
\text{Modus: }2 \text{ in } 3 \ (\text{bimodalno, } f=5)
\]
\[
\text{Mediana: }3
\]
\end{resitev}

\newpage

\begin{resitev}[Naloga 6 --- Kombinatorika in verjetnost]

V skupini je $6$ fantov in $4$ dekleta, skupaj $10$ oseb.

a) Vseh možnih odborov treh članov:
\[
\binom{10}{3}=120
\]
Odbori brez fantov (same dekleta):
\[
\binom{4}{3}=4
\]
Odbori z vsaj enim fantom:
\[
120-4=116
\]

b) Verjetnost, da so v odboru same dekleta:
\[
P=\frac{4}{120}=\frac{1}{30}\approx0{,}033
\]
\end{resitev}

\vspace{0.5cm}

\begin{resitev}[Naloga 7 --- Odvod in kot s koordinatno osjo]

a) $f(x)=4x^3-3x^2+2x$
\[
f'(x)=12x^2-6x+2
\]
\[
f'(1)=12-6+2=8
\]

b) $g(x)=\dfrac{4}{\sqrt[3]{x^2}}-3=4x^{-2/3}-3$
\[
g'(x)=4\cdot\left(-\frac{2}{3}\right)x^{-5/3}=-\frac{8}{3}x^{-5/3}
\]
Pri $x_0=-1$:
\[
(-1)^{-5/3}=\frac{1}{\left((-1)^{1/3}\right)^5}=\frac{1}{(-1)^5}=-1
\]
\[
g'(-1)=-\frac{8}{3}\cdot(-1)=\frac{8}{3}
\]
Kot med tangento in abscisno osjo:
\[
\varphi=\arctan\!\left(\frac{8}{3}\right)\approx 69{,}44^\circ
\]
\end{resitev}

\newpage

\begin{resitev}[Naloga 8 --- Enačba tangente]

$f(x)=x^2-1$, \quad $x_0=3$
\[
f(3)=8,\qquad f'(x)=2x,\qquad f'(3)=6
\]
Enačba tangente:
\[
t:\;y=6(x-3)+8=\boxed{6x-10}
\]

\begin{center}
\begin{tikzpicture}[scale=0.7]
\begin{axis}[
axis lines=center,
xlabel={$x$}, ylabel={$y$},
xmin=-1, xmax=5,
ymin=-3, ymax=20,
xtick={-1,0,...,5},
ytick={0,4,...,20},
grid=both,
grid style={line width=0.2pt, draw=gray!30},
width=10cm, height=8cm,
]
\addplot[domain=-0.5:4.5, samples=100, thick, color=brown, line width=1.5pt]
  {x^2-1} node[above left, pos=0.85] {$f(x)=x^2-1$};
\addplot[domain=1:4.5, samples=50, thick, color=blue, line width=1.5pt]
  {6*x-10} node[right, pos=0.9] {$t:y=6x-10$};
\addplot[only marks, mark=*, mark size=2.5pt, color=red]
  coordinates {(3,8)};
\node[above left, color=red] at (axis cs:3,8) {$(3,8)$};
\end{axis}
\end{tikzpicture}
\end{center}
\end{resitev}

\end{document}
